\section{When order matters: certificates in a free algebra}
\label{sec:noncommutative}

Every algebraic certificate so far has been commutative. Cone and quadratic
certificates multiply polynomials in $\Q[x_1,\ldots,x_n]$; the closure workers of
Section~\ref{sec:closure} multiply matrices, but the \emph{coefficients} they
manipulate commute, and the one place where order was forced to matter --- the
Ore rule $E_n n=(n+1)E_n$ of Section~\ref{sec:ore} --- was a single prescribed
commutation law between one shift and its coefficients.

That is a narrow escape from a general problem. Operators, matrices of unknown
dimension, quantum observables, and program effects all live in algebras where
$ab\ne ba$ and where no commutation law is given in advance. Two contributing
efforts take the general case: the free associative algebra $\Q\langle
X\rangle$ on a finite alphabet, presented by relations $r_1,\ldots,r_s$, with
the question
\[
 r_1=\cdots=r_s=0 \quad\Longrightarrow\quad p=0
\]
or, on the positivity side, whether $p$ is a sum of Hermitian squares modulo
those relations. Neither is a new area of mathematics --- the Diamond Lemma and
noncommutative Positivstellens\"atze are established. What is new here is that
both are cut down to a certificate a small checker can verify against the
\emph{original} relations.

\subsection{The two-sided certificate, and one theorem for two searches}

In a commutative ring an ideal member is $\sum_i u_i f_i$. In a free algebra it
is not: a relation may be multiplied on both sides, so the certificate carries
\emph{two-sided contexts},
\begin{equation}\label{eq:two-sided}
 p=\sum_{k} \alpha_k\, u_k\, r_{i_k}\, v_k,
\end{equation}
with words $u_k,v_k$, rational $\alpha_k$, and each $r_{i_k}$ one of the
\emph{supplied} relations. The checker expands \eqref{eq:two-sided} in the free
algebra and compares with $p$. It never needs a basis, a term order, a
confluence property, or a termination argument.

\begin{theorem}[Two-sided soundness]\label{thm:nc-sound}
If \eqref{eq:two-sided} holds in $\Q\langle X\rangle$ and every $r_i$ vanishes
in an algebra $A$ under a homomorphism fixing the generators, then $p$ vanishes
in $A$.
\end{theorem}
\begin{proof}
A homomorphism sends the right side of \eqref{eq:two-sided} to
$\sum_k\alpha_k\,\varphi(u_k)\,\varphi(r_{i_k})\,\varphi(v_k)$, and each
$\varphi(r_{i_k})=0$. Multiplication by $\varphi(u_k)$ on the left and
$\varphi(v_k)$ on the right is well defined in any associative algebra, so
every summand vanishes and $\varphi(p)=0$.
\end{proof}

Two of the contributing efforts prove this theorem independently, with the same
proof, and it is stated once here. What differs is entirely the \emph{search}
that produces \eqref{eq:two-sided}, and the two searches disagree about whether
completion is needed at all.

\paragraph{Lane one: eliminate contexts, do not complete.} \src{r2} fixes a
context bound, generates the finite set of products $u r_i v$ available within
it, and solves one exact linear system over the resulting word coefficients.
Its own statement of why this suffices is worth quoting, because it is a
deliberate refusal:

\begin{quote}\small
This proof does not appeal to termination of noncommutative completion or to a
Noetherian property. It isolates a finite-dimensional vector-space problem at
the exact required degree. It needs no global completion result for the
presented algebra.
\end{quote}

\paragraph{Lane two: complete, under a cap.} \src{r3} runs bounded,
target-oriented Gr\"obner--Shirshov completion, tracing each new rule back to
the original relations so that the emitted certificate is still in the form
\eqref{eq:two-sided}. It is equally explicit that the process it is running has
no termination guarantee:

\begin{quote}\small
The \emph{completion process} need not terminate. \ldots{} The prototype
therefore caps reduction steps, accumulated term/provenance size, rules, and
examined compositions.
\end{quote}

These are not rival implementations of one algorithm and the merge does not
pick a winner. A supplied orientation of the hypotheses can be terminating and
still incomplete for a given target, which is exactly when critical pairs expose
a consequence that no direct reduction of the target would reach --- lane two's
case. Equally, when all relations are homogeneous, the answer lives at one known
degree and completion is wasted work --- lane one's case. The routing rule is
the obvious one and one of the efforts already states it: if every relation is
homogeneous, solve at the target degree; otherwise try bounded contexts, and
fall back to capped completion.

\begin{remark}[Centrality is a hypothesis, not a formatting choice]
Both efforts arrive independently at the same warning, and it is the kind of
thing that is easy to lose in translation to Lean. The certificates use
rational coefficients, and \eqref{eq:two-sided} implicitly assumes those
rationals \emph{commute with everything}. In an arbitrary ring they need not.
The source bridge therefore has to establish that the target algebra is a
$\Q$-algebra with central rational scalars; ``the coefficients are numbers'' is
not an argument.
\end{remark}

\subsection{A refutation object where order matters}

A nonzero remainder is not a disproof: it means this search, at this bound,
found nothing. Both efforts say so, and \src{r3} labels its completion-disabled
control \unknown{} rather than \textsc{refuted} precisely to make the point.

For \emph{homogeneous} relations, though, \src{r2} gets something better. At a
fixed degree the consequences of the relations form a finite-dimensional
subspace, computable exactly, so the search is \emph{degree-complete}: if $p$ is
not in that subspace, no certificate of any size exists at that degree. The
refutation is then a concrete matrix representation --- the truncated quotient
algebra itself --- in which every relation vanishes and $p$ does not.

That is a genuine addition to the collection's stock of negative evidence. The
Kripke countermodels of Section~\ref{sec:structural} refute derivability in a
propositional calculus; a separating word refutes an automaton equivalence; this
refutes an algebra identity, by exhibiting an algebra. All three are checkable
by re-evaluation, and none of them is an absence of proof dressed up as a
disproof.

\subsection{Positivity when the variables do not commute}

The other half of the noncommutative problem is inequality rather than equality.
For a Hermitian $p$ of degree $2d$ over an alphabet of size $m$, the object that
replaces the commutative Gram matrix is indexed by \emph{words} of length at
most $d$ rather than by monomials:
\[
 p=\sum_{u,v} G_{uv}\, u^{*}v,
\]
and the essential difference is that this representation is \emph{unique}. In
the commutative case distinct monomial pairs collide --- $x\cdot y$ and $y\cdot
x$ contribute to the same coefficient --- and recovering a Gram matrix requires
solving for the freedom that collision creates. In the free algebra
$u^{*}v$ determines $(u,v)$, so the matrix is read straight off the
coefficients, over $m^d$ words with no collisions to resolve.

A certificate is then an exact rational $LDL^{\mathsf T}$ decomposition of that
matrix with nonnegative diagonal, which the checker re-multiplies. Two
extensions matter for real targets. Constrained positivity uses adjoint
sandwiches $q^{*}gq$ rather than plain squares, so a supplied positive element
$g$ can be used in the middle. And when the alphabet carries an involution,
the generator of the relevant cone is a structured ray rather than an arbitrary
square, which keeps the search finite.

\begin{remark}[Dimension-free is the point]
These certificates are valid \emph{independently of the dimension} of any
matrix representation. That is the reason to work in the free algebra at all,
and it is what a numerical semidefinite solver at one fixed size cannot give.
The recorded corpus includes a Bell-inequality (CHSH) derivation, where exactly
this matters: the bound is wanted for every representation, not for one chosen
Hilbert space.
\end{remark}

\subsection{What these workers do not do}

The scope conditions are narrower than the machinery suggests, and the efforts
state them.

Neither lane refutes anything in the inhomogeneous case. A bound was reached; a
larger bound may succeed; and the \unknown{} verdict carries no information
about the goal's truth.

Degree-completeness is completeness \emph{at a degree}, not completeness. A
homogeneous target that needs consequences at a higher degree is outside the
guarantee, and the guarantee says nothing about the inhomogeneous fragment at
all.

An $LDL^{\mathsf T}$ certificate establishes that $p$ is a sum of Hermitian
squares in the free algebra. Transporting that to a statement about a particular
operator algebra needs the representation, the involution and the positivity
order to match the source, and a trace-positivity claim needs more than a
sum-of-squares decomposition. These are separate obligations and the article
does not fold them into one.

Finally, and as everywhere in this document: none of this has been checked by
Lean. Mathlib has a free algebra, and the two efforts cite it under two
different module paths without either having compiled against it --- a small but
exact illustration of why the pinned-revision discipline of
Section~\ref{sec:trust} exists.
